\documentclass[english, a4paper]{article}
\usepackage[T1]{fontenc}
\usepackage[latin9]{inputenc}
\usepackage{geometry}
\usepackage{url}
\usepackage{babel}
\usepackage{csquotes}
\usepackage{longtable}
\usepackage{booktabs}
\usepackage[]{biblatex}
\addbibresource{bibliography.bib}
\geometry{verbose,tmargin=3cm,bmargin=3cm,lmargin=3cm,rmargin=3cm}

\title{
    Research Plan for CSE3000 Research Project\\
    {\Large \textit{Reducing data in visual AI: evaluating the data efficiency of Barlow Twins}}
}
\author{Yan Olerinskiy}
\date{April 2026}

\begin{document}
\maketitle


\section*{Background of the research}
Modern visual AI is powered by very large datasets, but access to such datasets is highly uneven: most web-scale image data sits with a handful of large companies, while universities and independent researchers have far more limited means to train competitive models. Large datasets are also hard to curate, which complicates efforts to control for fairness, copyright, and bias.

This project investigates how to train visual AI more data-efficiently. The common recipe is a Vision Transformer (ViT) backbone pre-trained in a self-supervised way on a generic dataset, and then adapted to a downstream task with Parameter-Efficient Fine-Tuning (PEFT). Existing self-supervised methods have mostly been evaluated at scale, and have not been carefully compared in the small-data regime that universities actually live in. Following the ``scaling down deep learning'' idea, we deliberately reduce both the data and the compute, and produce learning curves (downstream score vs. pre-training dataset size) for several popular self-supervised methods under matched conditions.

\section*{Research Question}
\textbf{Main group question.} How can the data efficiency of small-compute, self-supervised pre-trained visual foundation models be evaluated using learning curves on a set of downstream tasks?

\textbf{Individual sub-question.} How does the data efficiency of Barlow Twins compare to the alternative self-supervised methods studied by the rest of the group (DINO, MAE, MoCo v3, I-JEPA), when used to pre-train a ViT-Tiny/8 backbone on Tiny ImageNet and adapted to CIFAR-10 with PEFT?

\textbf{Why this is achievable in the timeframe.} The scope is realistic for a single quarter because we deliberately keep things small: Tiny ImageNet rather than full ImageNet, ViT-Tiny/8 rather than larger ViTs, and CIFAR-10 (with VTAB-style additions only if compute permits) rather than the full VTAB benchmark. We reuse existing implementations from Hugging Face, the \texttt{solo-learn} library, and the original Barlow Twins codebase rather than building from scratch, and the five group members share one data loader and one downstream evaluation script, so my individual workload is essentially one method, several pre-training fractions, one downstream task.

\textbf{Success criteria.} The minimum viable result is a learning curve for Barlow Twins with pre-training dataset fraction on the $x$-axis and CIFAR-10 PEFT accuracy on the $y$-axis. The ideal final result is the group-level aggregated plot, in which my Barlow Twins curve is overlaid on the curves of the other four methods, so that we can read off the data regimes in which Barlow Twins is more or less data-efficient than the alternatives.

\section*{Method}
The work is split into a shared group pipeline and individual method-specific experiments.

\textbf{Shared pipeline (group task).} Together with the four other students (Leonid -- DINO, Makar -- MoCo v3, Maksim -- I-JEPA, Dimo -- MAE), I will help build a shared evaluation harness so that all five sets of results are directly comparable. The harness consists of a PyTorch data loader that takes a fixed seed and slices Tiny ImageNet into agreed fractions (e.g.\ 5\%, 10\%, 25\%, 50\%, 100\%) so that everyone pre-trains on the exact same image subsets; a downstream PEFT fine-tuning script that takes a pre-trained ViT-Tiny/8 backbone and reports CIFAR-10 accuracy under a fixed protocol; and a simple results format (CSV with method, fraction, seed, accuracy) that is easy to merge into the final aggregated plot.

\textbf{Individual work on Barlow Twins.} On top of the shared pipeline, my individual contributions are: set up Barlow Twins on the shared ViT-Tiny/8 backbone, starting from an existing reference implementation (likely \texttt{solo-learn}) and adapting it to the shared data loader; pre-train Barlow Twins on each of the agreed Tiny ImageNet fractions, keeping a single sensible hyperparameter configuration (per supervisor advice, no broad sweep) and using early stopping when the SSL loss flattens; fine-tune each pre-trained checkpoint on CIFAR-10 through the shared PEFT script and record the test accuracy; plot my individual Barlow Twins learning curve and submit the numbers in the shared CSV format so the group can produce the aggregated comparison; and, if time permits, run a short extension experiment looking at how Barlow Twins's projection-head dimension and decorrelation strength interact with reduced data.

\textbf{Tools and data.} PyTorch as the main framework; Hugging Face \texttt{transformers} and \texttt{peft} for the backbone and the PEFT method; \texttt{solo-learn} as a starting point for the Barlow Twins ViT recipe; Tiny ImageNet (\texttt{zh-plus/tiny-imagenet}, 64$\times$64) for pre-training; CIFAR-10 (a VTAB-natural task) upsampled to 64$\times$64 for downstream evaluation. I will start on a local machine and only ask the supervisors for cluster access if it becomes a real bottleneck.

\section*{Planning of the research project}
The plan covers the 10 weeks of Q4. The internal target is to have the group's aggregated learning-curve plot ready by the end of Week~6, so that Weeks~7--10 can be spent on the writing-heavy course deliverables and, if time allows, on optional improvement experiments. Course deadlines are taken from the Brightspace schedule for our cohort.

{\small
\begin{longtable}{p{1.2cm} p{9.5cm} p{3.5cm}}
\toprule
\textbf{Week} & \textbf{Activity / Deliverable} & \textbf{Category} \\
\midrule
\endhead

Week 1 & Survey self-supervised methods and reference implementations to choose one to claim & Project Task \\
       & Peer meeting: claim method, agree on the shared stack (backbone, dataset, fractions) & Peer Meeting \\
       & Start contributing to the shared data loader & Project Task \\
       & Supervisor meeting & Supervisor Meeting \\
       & Submit Research Plan & Course Deadline \\
\midrule

Week 2 & Wrap up the shared pipeline and get Barlow Twins running end-to-end at one fraction & Project Task \\
       & Read data-efficiency literature in parallel (e.g.\ ``Scaling Down Deep Learning'' and follow-ups) & Project Task \\
       & Supervisor meeting -- discuss MVP scope and main risks & Supervisor Meeting \\
       & ACS Assignment 1; Research Plan Presentation & Course Deadline \\
\midrule

Week 3 & Pre-train Barlow Twins across the agreed Tiny ImageNet fractions; iterate as issues arise & Project Task \\
       & Supervisor meeting & Supervisor Meeting \\
       & ACS Assignments 2a, 2b & Course Deadline \\
\midrule

Week 4 & Continue pre-training; investigate any unexpected behaviour & Project Task \\
       & Peer meeting: cross-check that intermediate numbers across methods are comparable & Peer Meeting \\
       & Supervisor meeting; confirm Final Presentation date & Supervisor Meeting \\
       & Session ACS (Poster) & Course Deadline \\
\midrule

Week 5 & Run downstream PEFT evaluation; produce my Barlow Twins learning curve and midterm slides & Project Task \\
       & Supervisor meeting -- dry-run the midterm narrative & Supervisor Meeting \\
       & Midterm Presentation; process feedback; ACS Assignment 3 & Course Deadline \\
\midrule

Week 6 & Peer meeting: merge all five learning curves into the aggregated plot & Peer Meeting \\
       & \textbf{MILESTONE: aggregated learning-curve graph (baseline pass)} & Project Milestone \\
       & Supervisor meeting -- discuss findings and direction for the paper & Supervisor Meeting \\
\midrule

Week 7 & Write Paper Draft v1 (Intro, Method, Baseline Results) & Course Task \\
       & (Optional) Begin exploring extensions: projection-head dimension, decorrelation strength, augmentation choices & Advanced Research \\
       & Coaching session; supervisor meeting & Meetings \\
       & Submit Paper Draft v1; peer-review another group's draft; Feedback v1 review & Course Deadline \\
\midrule

Week 8 & Integrate peer-review feedback; write Discussion, Conclusion, and limitations & Course Task \\
       & Coaching session; supervisor meeting & Meetings \\
       & Submit Paper Draft v2 & Course Deadline \\
\midrule

Week 9 & Address Responsible Professor feedback on v2; polish figures and reproducibility appendix; sketch poster & Project Task \\
       & Coaching session; supervisor meeting & Meetings \\
       & Feedback v2 review (Responsible Professor); Final Paper submission & Course Deadline \\
\midrule

Week 10 & Finalise poster; rehearse Final Presentation & Project Task \\
        & Poster Submission; Final Presentation for the supervising team & Course Deadline \\

\bottomrule
\end{longtable}
}

\printbibliography

\end{document}